\section{Discussion}
\label{sec:implications}

\subsection{Implications for Mechanistic Interpretability}

Our findings have several implications for the field:

\paragraph{Superposition is fundamental, not artifactual.} The persistence of superposition across toy autoencoders, transformers, translation models, and continuous-thought architectures (Table~\ref{tab:full_comparison}) confirms it is a general strategy neural networks employ to manage limited capacity, not an artifact of simplified experimental settings. Every model with $n > d$ exhibits non-trivial superposition, with interference structure governed by the compression ratio.

\paragraph{Sparsity drives superposition.} In all architectures, increased sparsity leads to more superposition. The phase diagram (Figure~\ref{fig:phase_diagram}) shows a clear transition at $S \approx 0.9$, consistent with the critical sparsity threshold of Corollary~\ref{cor:capacity}. This aligns with the theoretical prediction that interference costs decrease when features are rarely co-active.

\paragraph{Practical models use superposition.} The translation model's ability to maintain performance through a 32$\times$ compression (512$\rightarrow$16 dimensions) demonstrates that superposition enables efficient encoding of real linguistic features, not just synthetic ones. The quantitative analysis further reveals that 85.4\% of feature pairs remain near-orthogonal even at 2$\times$ compression in the 256-dimensional bottleneck variant.

\paragraph{Transfer of techniques.} The equivalence theorem (Theorem~\ref{thm:equivalence}) and the shared geometric structure observed in Section~\ref{sec:empirical} suggest that methods developed for representational superposition may apply to reasoning superposition. Sparse autoencoders \citep{bricken2023monosemanticity} disentangle superposed features by learning an overcomplete basis; analogous methods could disentangle superposed reasoning traces in continuous-thought models. Feature steering \citep{templeton2024scaling} modifies activations to change model behavior; similar interventions on thought vectors could steer reasoning.

\subsection{Implications for AI Safety}

Understanding superposition is crucial for AI safety efforts. First, superposition complicates the identification of individual features; techniques like sparse autoencoders \citep{bricken2023monosemanticity} aim to disentangle superposed representations. Second, \citet{templeton2024scaling} showed that identified features can be used to steer model behavior, and our results suggest this approach should generalize across architectures. Third, the capacity bounds (Corollary~\ref{cor:capacity}) predict when interference-induced errors occur: when the number of active objects exceeds $d + d^2 S / \tau$, superposition causes errors. This could enable proactive safety measures.

If language models reason via superposition of traces \citep{zhu2025reasoning}, then harmful reasoning may be one of multiple superposed traces. Steering away from harm requires identifying and suppressing specific traces without disrupting beneficial ones. The co-occurrence structure $P$ determines which traces can be independently manipulated---a direction our equivalence framework makes precise.

\subsection{Implications for Architecture Design}

\paragraph{Width vs.\ depth tradeoff.} Theorem~\ref{thm:timespace} and our empirical results (Section~\ref{sec:timespace_empirical}) quantify a fundamental tradeoff. Wider models ($d$ large) require less superposition and achieve lower interference per feature. Deeper or iterative models ($T$ large) can exponentially amplify capacity through gating, achieving interference comparable to models with $2$--$4\times$ the bottleneck dimension. For interpretability, wider models may be preferable (less superposition to disentangle), but iterative models may be more parameter-efficient.

\paragraph{Bottleneck design.} Our results suggest deliberate bottlenecks can induce structured, analyzable representations. The bottleneck forces superposition by constraining $d < n$, and the resulting weight matrices exhibit systematic geometric organization (near-equiangular frames at high sparsity, POS-clustered embeddings in translation). The bottleneck thus acts as an ``interpretability checkpoint'' where all information must pass through a low-dimensional space amenable to analysis.

\paragraph{Practical guidance.} The Welch ratio data (Table~\ref{tab:interference_metrics}) indicates that learned representations operate at 4--8$\times$ the theoretical minimum coherence. Architectures that approach tighter packing---through architectural inductive biases or auxiliary losses penalizing $\mathcal{L}_{\text{int}}$---could improve both capacity and interpretability simultaneously.
